\chapter{Introduction}
\label{intro.ch}

Digital electronics increasingly occupy safety- and mission-critical roles in radiation
environments. Nuclear energy systems rely on embedded computation for sensing, data
acquisition, actuator control, and condition monitoring. Space and high-altitude platforms
similarly depend on autonomous processing to maintain communications, navigation, and
health management under particle radiation. In these settings, radiation-induced faults are
not merely a reliability nuisance; they can compromise observability of the physical
process, corrupt stored state, or drive a controller into an unsafe or unknown operating
condition. As modern systems move toward higher integration and tighter coupling between
software-defined functionality and hardware, the ability to characterize and mitigate
radiation effects at the \emph{system} level becomes increasingly
important~\cite{baumann2005radiation, schwank2013radiation}.

Field-programmable gate arrays (FPGAs) are a compelling technology for harsh-environment
instrumentation and control because they combine high-throughput digital processing with
reconfigurability. A single FPGA can integrate sensor interfaces, deterministic timing,
data reduction, fault detection, and control logic that would otherwise require multiple
discrete components. However, many widely used FPGA families are susceptible to
radiation-induced single-event effects (SEEs), including transient upsets in sequential
logic and persistent upsets in configuration
memory~\cite{kastensmidt2014survey, ostler2009sram}. These effects can manifest as
straightforward bit flips in data paths, but they can also appear as subtle and difficult-to-
diagnose failures such as stalled state machines, corrupted mode registers, or sporadic
timing violations triggered by transients. Consequently, evaluating FPGA-based systems
using only simplistic benchmarks or post-run inspection can miss the failure mechanisms
that matter most in real deployments~\cite{artola2021benchmarks}.

This dissertation is motivated by a practical question: \emph{how can FPGA-based digital
systems be evaluated under irradiation in a way that preserves in-situ operation, provides
high observability of faults and performance degradation, and yields data products that
support quantitative analysis?} Addressing this question requires more than a single
measurement technique. It demands an experimental workflow that spans on-chip
instrumentation, robust telemetry, test automation, and analysis methods that connect
radiation events to system-level outcomes. The remainder of this chapter provides the
system context for the problem, describes the radiation environments of interest, motivates
the use of FPGAs while acknowledging their vulnerabilities, and outlines the in-situ test
and monitoring philosophy that underpins the work.

\section{Radiation-tolerant instrumentation and control}
Digital instrumentation and control (I\&C) systems form the nervous system of complex
engineering facilities~\cite{nrc1997digital, iaea2016design}. In nuclear applications, I\&C platforms acquire measurements from
distributed sensors, estimate plant state, and command actuators that regulate power,
temperature, and flow. Even when safety systems are implemented with conservative
redundancy and independent protection channels, auxiliary digital subsystems are often
responsible for data logging, operator interfaces, and supervisory
control~\cite{korsah1999environmental}. In research and
test facilities, digital platforms also play a central role in experiment automation, beamline
instrumentation, and synchronized acquisition of diagnostic signals.

Radiation environments challenge I\&C systems in two distinct ways. The first is direct
exposure to particles and secondary radiation that can produce faults in electronics. The
second is the operational constraint that systems must often remain functional while data
are collected. In many experiments, the most valuable information occurs \emph{during}
exposure: transient anomalies, intermittent recoveries, and performance drift with
accumulated fluence. These behaviors are easy to lose when experiments rely on coarse
snapshots, infrequent configuration readback, or manual intervention.

For embedded digital systems, the most consequential failures are frequently those that
degrade correctness without producing an obvious halt. A sensor interface that silently
drops samples, a communications link that intermittently corrupts frames, or a controller
that remains stable but exhibits increased steady-state error can each undermine mission
objectives while evading simple ``pass/fail'' criteria. Therefore, a central requirement for
harsh-environment evaluation is \emph{observability}: the ability to measure internal state,
detect anomalies, and attribute system-level behaviors to underlying fault mechanisms.

In-situ observability must also be engineered within practical constraints. Telemetry
bandwidth is finite, particularly when long cables, optical links, or radiation-hardened
interfaces are required. Logging must tolerate resets and interruptions, and it must preserve
sufficient timing context to correlate faults with external conditions such as beam
transitions, actuator commands, or power supply events. These realities motivate a design
philosophy in which monitoring and control are co-designed with the experiment itself.

\section{Radiation environments and single-event effects in electronics}
Radiation fields relevant to terrestrial nuclear facilities, beamline experiments, and
aerospace missions vary widely in particle species and energy spectra. Fast neutrons are of
particular interest for terrestrial testing because they can produce displacement damage and
ionization through secondary charged particles generated in nuclear interactions with
silicon and package materials~\cite{normand1996single, ziegler1996terrestrial}. High-energy protons and heavy ions are often emphasized
in space applications because they can directly deposit charge along a track in sensitive
volumes, leading to transient and sometimes destructive events. Gamma radiation primarily
contributes to total ionizing dose (TID), which can shift device parameters and increase
leakage currents over time. In practical environments, mixed fields can be present, and
secondary particles produced in surrounding structures can dominate the effective spectrum
seen by the electronics.

Single-event effects are a family of phenomena in which a single particle interaction
perturbs device behavior~\cite{dodd2003basic, messenger1992effects}. In digital logic, the most commonly discussed SEE is the
single-event upset (SEU), in which a storage element such as a flip-flop or memory bit
changes state. When an upset occurs in a data register, it may appear as an incorrect numeric
value that propagates through computation. When it occurs in a state register, it can change
the sequencing of a controller, alter a mode selection, or prevent a process from advancing.
Single-event transients (SETs) occur when a particle strike produces a short-duration pulse
in combinational logic. These pulses may be filtered by logic depth and timing, or they may
be captured by sequential elements depending on clock phase and pulse
width~\cite{dodd2004production}. More severe
events such as single-event latchup (SEL) can induce a high-current state that persists until
power is cycled or protective circuitry intervenes, potentially damaging the device.

A defining challenge of SEE characterization is that observed behavior depends strongly on
system context. A single bit flip in a redundant arithmetic datapath may have little impact,
while an identical flip in a control enable or address pointer can halt a test or invalidate
measurements. Additionally, some failures are \emph{silent}: a system may continue to run
and produce outputs that appear reasonable, yet internal fault accumulation or altered timing
causes gradual performance degradation. These realities motivate system-level experiments
that measure not only event counts, but also the consequences of events on functional
behavior~\cite{natacad2018testing, keller2018dynamic}.

Quantitative analysis typically describes SEE sensitivity using an upset rate or cross-section
as a function of fluence and particle spectrum~\cite{reed2003ground, shaneyfelt2008qualification}. However, a single rate does not fully capture
the engineering question of interest: \emph{what does the system do after an event occurs?}
For embedded control and monitoring platforms, the post-event trajectory---recovery,
persistence, amplification, or self-correction---is often more important than the event count
itself. This dissertation therefore emphasizes measurement strategies that preserve and
interpret the temporal evolution of the system during irradiation.

\section{Field-programmable gate arrays as harsh-environment digital platforms}
FPGAs implement digital circuits using a fabric of programmable logic blocks, routing
resources, and embedded hard macros such as block RAM, DSP slices, and high-speed
serial transceivers. Their defining feature is reconfigurability: functionality is specified by
a configuration bitstream that programs lookup tables, interconnect, and optional features
of the device. This capability enables rapid iteration and consolidation of complex digital
subsystems, which is particularly attractive in experimental environments where
requirements evolve.

For instrumentation and control, FPGAs provide several practical advantages. Deterministic
timing can be achieved by implementing synchronous datapaths with precisely controlled
latency. Multiple heterogeneous functions---such as sensor sampling, filtering, state
estimation, actuator drive, and communications---can be integrated while maintaining
tight timing alignment. Hardware-level parallelism also supports high-rate data reduction
and event detection, enabling long-duration experiments without saturating telemetry
bandwidth.

These advantages come with a fundamental vulnerability: in many FPGA families, the
configuration memory that defines the circuit is implemented using radiation-sensitive
storage~\cite{kastensmidt2014survey, rakhman2017sram}. Upsets in configuration bits can alter logic functions or routing, creating persistent
errors that remain until the configuration is corrected. Upsets in user logic storage elements
can corrupt state, and SETs can perturb combinational paths. Consequently, the relevant
question is not whether an FPGA experiences upsets, but how the overall system is
architected to \emph{observe}, \emph{tolerate}, and when necessary \emph{recover} from them.

Mitigation strategies exist across multiple layers. At the circuit level, redundancy such as
duplication or triplication can mask certain classes of errors, though at the cost of resource
usage and potentially increased sensitivity to common-mode
faults~\cite{lyons1962tmr, sterpone2005tmr, pratt2008partial}. At the configuration
level, readback and scrubbing can detect and correct configuration upsets, but they require
additional infrastructure and can be difficult to integrate with strict timing
requirements~\cite{heiner2009scrubbing, berg2008scrubbing, stoddard2017hybrid}.
At the system level, watchdogs and health monitors can detect stalls or inconsistent outputs
and trigger controlled recovery actions. The selection of mitigations is inherently tied to
application constraints: a control loop may prioritize bounded recovery time, while a data
logger may prioritize data integrity and provenance.

A key observation motivating this dissertation is that many experiment-relevant failure
modes are \emph{control-path} failures rather than straightforward datapath bit flips. Control
paths include counters, state machines, scheduling flags, communication framing, and mode
registers. Upsets in these elements can prevent valid comparisons, stop scanning, or lock a
system into a non-responsive state without producing an unambiguous error signature in
the primary outputs~\cite{morgan2005seu}. Capturing these failure modes requires on-chip observability and
telemetry that are designed from the outset to detect ``stuck'' or ``silent'' conditions.

\section{In-situ monitoring and irradiation testing methodology}
In-situ irradiation testing aims to measure system behavior while the system operates under
exposure, rather than treating radiation as an off-line stress applied prior to a functional
check. This approach is essential when the effects of interest are transient, intermittent, or
dependent on accumulated fluence. It is also essential when recovery actions---such as
scrubbing, resets, or power-cycling after latchup---are part of the operational concept and
therefore must be measured as part of system performance.

A practical in-situ methodology must provide fault visibility that is commensurate with the
complexity of the system. For FPGA platforms, this motivates instrumentation that can
detect and summarize fault activity within large collections of sequential elements and
internal control structures without requiring full-state readout. The methodology must also
provide time-aligned telemetry so that fault indicators can be correlated with external
conditions, control commands, and measured plant response. Finally, it must support
repeatable experimentation, including clear run segmentation, metadata capture, and
post-processing workflows that reduce ambiguity.

One approach to fault visibility is to embed structured test logic within the FPGA design.
Large sets of flip-flops can be organized into scan structures or signature-based monitors,
enabling efficient detection of deviations from expected behavior. Complementary watchdog
logic can supervise progress through known operational phases, flagging stalls or invalid
states that might otherwise appear as a simple cessation of log messages. Telemetry formats
can be designed to favor robustness and interpretability, prioritizing concise, periodic
reports that survive occasional corruption and support automated parsing.

While fault detection infrastructure can quantify upset activity, system relevance is often
best demonstrated using a meaningful workload. In this dissertation, closed-loop control is
used as a representative workload because it couples internal computation to externally
observable behavior. A controller operating on an FPGA interacts with a plant model or
physical system through measured feedback. Radiation-induced faults can therefore be
evaluated not only by their occurrence, but by their impact on performance metrics such as
tracking error, overshoot, actuator saturation, and stability margins. Importantly, a closed-
loop workload can reveal fault masking and fault amplification: some errors are corrected
by feedback, while others are reinforced by the loop dynamics.

The combination of on-chip fault monitors, watchdog-based health supervision, robust
telemetry, and a closed-loop workload provides a foundation for system-level irradiation
studies that move beyond binary outcomes. Rather than asking whether an FPGA ``survived''
a run, the experiment can ask how performance evolves with exposure, which subsystems
dominate observed failures, and which mitigations provide the most impact per unit
resource cost.

\section{Scope and organization of this dissertation}
This dissertation focuses on radiation-induced effects in FPGA-based digital systems, with
an emphasis on in-situ experimentation and system-level interpretation of faults. The work
considers single-event phenomena that manifest during operation and the architectural
choices that influence detectability, tolerance, and recovery. While total ionizing dose and
displacement damage can influence long-term behavior, the primary emphasis is placed on
single-event-driven disruptions relevant to real-time operation.

Chapter~2 provides background on the physics of radiation interactions relevant to digital
electronics and reviews prior work in FPGA radiation testing and mitigation. Chapter~3
describes the hardware and software infrastructure used to conduct in-situ experiments,
including embedded monitors and telemetry workflows. Chapter~4 presents the closed-loop
FPGA workload used to connect internal fault activity to externally measurable performance.
Chapter~5 presents irradiation results and analysis, including event classification and
quantitative metrics versus exposure. Chapter~6 concludes the dissertation with a synthesis
of findings and directions for future work in harsh-environment digital platforms and
radiation-aware system design.